\section{Open Challenges and Future Directions}

The review and Wuhan benchmark point to several challenges that are now becoming central for satellite embeddings. The first is temporal ambiguity. Many current systems use the same word, embedding, for annual products, short seasonal composites, single-date observations, and multi-step time series. These objects can all be useful, but they do not preserve the same temporal information. A representation that averages a year of observations may be stable for land-cover mapping but poorly suited to flood response or phenology-sensitive agriculture; a temporal encoder may preserve seasonal dynamics but require inputs that are difficult to obtain in cloudy or data-sparse regions; a precomputed product may be easy to sample but locked to the upstream production calendar. Future papers should therefore report not only the sensor and spatial resolution of an embedding, but also its temporal support, compositing rule, update frequency, and intended interpretation. This would make model selection more transparent and would prevent annual products, scene encoders, and temporal sequence models from being compared as if they were the same representation.

The second challenge is scale mismatch. Remote sensing reviews have long emphasized that spatial, spectral, and temporal resolution determine what can be observed, but foundation-model discussions sometimes hide this fact behind architecture names. The Wuhan benchmark makes the issue concrete: embeddings produced at 10\,m, 30\,m, 500\,m, 1\,km, or 0.25$^\circ$ are not interchangeable, and resampling cannot create local information that was never present in the native representation. This is especially important for precomputed products because their convenience can encourage users to sample them outside the scale at which they are meaningful. A stronger evaluation culture would report native support, resampling choices, and target grid size together with accuracy. It would also include visual diagnostics, because map coherence often reveals scale failure before a single aggregate metric does.

The third challenge is the comparison between products and checkpoints. Precomputed embedding fields optimize access, coverage, and repeatability; on-the-fly encoders optimize flexibility, temporal control, and adaptation to new inputs. Both are legitimate satellite embedding regimes, but they should not be evaluated only by the same predictive score. Product models should be credited for low inference burden, stable large-area sampling, and reproducible feature layers. Checkpoint models should be credited for sensor-specific control, temporal customization, dense readout options, and the ability to adapt to new data streams. Future benchmarks should therefore pair quality metrics with deployment metrics such as extraction latency, GPU memory, missing-data behavior, storage footprint, and product availability. This is not an engineering afterthought. For Earth observation, these constraints often decide whether a representation can be used at continental scale or in repeated monitoring.

The fourth challenge is application realism. Standard remote sensing datasets remain valuable, but satellite embeddings are increasingly used in settings that do not look like benchmark leaderboards: sparse labels, mixed pixels, shifting seasons, changing sensors, region transfer, and continuous environmental quantities. The Wuhan benchmark is only a first step, yet it shows why related task families are useful. Classification, regression, water extraction, and segmentation expose different failure modes over the same region. Future work should extend this logic to larger and more diverse ROIs, cross-year splits, few-label protocols, and physically meaningful regression targets such as canopy height, biomass, crop condition, surface water dynamics, and urban intensity. The field also needs clearer reporting of negative results. A model that fails because of scale mismatch, unavailable inputs, or temporal misalignment teaches a different lesson from a model that fails because its representation is weak.

Finally, satellite embeddings need more standardized reporting. A useful model card for this area should include the pretraining data regime, sensor assumptions, native spatial and temporal support, output forms, embedding dimensionality, dense-readout availability, licensing or product access, and recommended downstream usage. Recent systems such as Prithvi-EO-2.0 and FlexiMo show that the field is continuing to expand toward more flexible multitemporal and multimodal representations \citep{prithvieo2_tgrs,fleximo_tgrs}. As that expansion continues, the review problem will become less about whether a model is a foundation model and more about whether its embedding is the right geospatial object for a given task. The next generation of surveys should therefore combine taxonomy, evidence, and reporting standards rather than treating model inventories as sufficient.
